\begin{table}[H]
\centering
\small\setstretch{1.0}
\caption{Self-Reported Racial Identity of Adults of Asian Ancestry by Generation \label{tab:asianbygen-adults}}
\begin{threeparttable}
\begin{tabular}{lccccc}
\toprule
 & & \multicolumn{4}{c}{Share self-reporting each identity} \\
\cmidrule(l{3pt}r{3pt}){3-6}
 & Observations & Asian only & White only & \specialcell{Asian and \\ White} & Other\\
\midrule
\textbf{First generation} & 286,553 & 0.95 & 0.03 & 0.00 & 0.02\\
\textbf{Second generation} & 102,372 & 0.70 & 0.15 & 0.10 & 0.05\\
\hspace{1em}Both parents Asian-born & 62,815 & 0.93 & 0.03 & 0.01 & 0.03\\
\hspace{1em}One Asian-born parent & 39,557 & 0.33 & 0.34 & 0.24 & 0.09\\
\hspace{1em}\hspace{1em}Asian father--White mother & 14,289 & 0.36 & 0.36 & 0.18 & 0.10\\
\hspace{1em}\hspace{1em}White father--Asian mother & 25,268 & 0.32 & 0.33 & 0.27 & 0.09\\
\bottomrule
\end{tabular}
\begin{tablenotes}
\small
\item[1] The samples include people of Asian ancestry ages 18 and above. First-generation Asian immigrants were born in an Asian country. Native-born second-generation Asian Americans have at least one parent born in an Asian country; ``One Asian-born parent'' rows require the other parent to be born outside Asia (the AW and WA labels follow the paper's convention). Grandparents' birthplaces are not observed for adults, so the third generation cannot be identified.
\item[2] Columns report the share of each row's observations in the racial identity categories used in the multinomial logit: ``Asian only,'' ``White only,'' and ``Asian and White'' (any combination of Asian and White, with or without Pacific Islander). ``Other'' collects Asian or Pacific Islander, other single races, and multiracial combinations without both an Asian and a White component; the multinomial logit excludes this category. Shares sum to one across the four columns.
\item[3] Data source is the 2004--2021 Current Population Survey.
\end{tablenotes}
\end{threeparttable}
\end{table}
